\chapter{Tension de surface}

\section{Origine de la tension de surface}

\section{Point de vue macroscopique}

À l'intérieur d'un fluide, comme nous l'avons vu, les \cindex{interactions} entre les atomes ou les molécules se manifestent, à l'échelle macroscopique, par de la \emph{viscosité}. Dans les fluides courants tels que l'air, l'eau ou l'huile, on peut admettre que la viscosité est isotrope et linéaire : c'est l'hypothèse newtonienne. À l'échelle microscopique, les interactions résultent de la combinaison de forces répulsives, à courte distance et associées à l'interpénétration des couches électroniques, et de forces attractives telles que les forces de Van der Waals et les ponts hydrogène, qui agissent à plus grande distance.

De façon un peu schématique, nous pouvons considérer que les forces attractives se manifestent la plupart du temps, alors que les forces répulsives se manifestent lorsque les molécules s'entrechoquent. On pourrait comparer ces dernières aux rebonds élastiques entre des billes qui se percutent.

Adoptons maintenant un point de vue plus macroscopique. Nous savons que si nous plaçons une surface témoin à l'intérieur du fluide, cette dernière va subir une force proportionnelle à la surface exposée. En grande partie, cette force exprime l'échange de quantité de mouvement causé par les chocs des molécules. Il s'agit de la pression cinétique, notée $p_k$, que l'on peut estimer en physique statistique par $p_k = \rho k_B T$ (où $k_B$ est la constante de Boltzmann).

Dans les liquides, cette pression est modifiée par les interactions intermoléculaires dominées par des phénomènes d'attraction. On peut donc écrire \cite{berry71aa} :
\begin{equation}
  p = p_k + p_f,
  \label{eq:pression}
\end{equation}
où $p_f$ est la pression dite "statique", associée aux interactions intermoléculaires d'attraction.

Considérons une surface témoin : la pression "statique" résulte de la force exercée par les molécules d'un côté de la surface (par exemple, à gauche) sur les molécules de l'autre côté. Cette pression statique est généralement négative (car les forces sont attractives). Ce n’est qu’en cas de pressions extrêmes, où la distance intermoléculaire est fortement réduite, que la pression statique peut devenir positive.

À l'intérieur du fluide, loin de la surface ou de l'interface avec d'autres matériaux, nous pouvons faire l'hypothèse que le milieu est \cindex{isotrope}. C'est d'ailleurs pour cette raison que les forces exercées contre une surface témoin sont indépendantes de l'orientation de cette surface. Près de la surface, l'isotropie est rompue. Il faut donc définir une pression "normale" $p_n$, qui s'exerce sur une surface témoin perpendiculaire à l'interface, et la distinguer d'une pression "tangentielle" $p_t$.

\section{Point de vue macroscopique : pression statique et dynamique}

Le point de vue macroscopique nous amène à réfléchir à la distribution de la pression. L’hypothèse centrale est que la pression cinétique $p_k$ dépend de la densité du fluide. Or, cette densité diminue près de la surface, en raison de l’équilibre dynamique entre les phases liquide et vapeur. Nous supposons également que $p_k$, qui résulte des chocs moléculaires, reste isotrope.

Nous avons déduit que $p_k$ diminue près de la surface. Par ailleurs, le fluide étant au repos, il n’y a pas d’accélération verticale macroscopique. Par conséquent, la dérivée partielle de la pression normale par rapport à $z$ est nulle : $\partial p_n / \partial z = 0$. Comme $p_k$ diminue près de la surface, $p_{f,n}$, la composante normale de la pression statique, doit être moins négative, donc augmenter en valeur absolue.

Cela est cohérent avec la figure 7 de \cite{marchand11aa} (REMETTRE CETTE FIGURE) : si on considère une surface témoin horizontale, les attractions étant à longue portée, il y a moins de molécules à "tirer" au fur et à mesure qu'on se rapproche de la surface. Autrement dit, si l'on s'intéresse à la force exercée sur la surface supérieure, il y a moins de molécules qui tirent vers le haut.

En revanche, la symétrie des forces d'attraction est préservée dans la direction tangentielle : $p_{f,t}$ n'a aucune raison de varier quand on s'approche de la surface. Donc, $p_k$ diminue (les forces "répulsives"), mais $p_{f,t}$ reste constant : cela revient à dire qu'il y a une force tangentielle nette qui s'exerce sur un élément de volume en surface, dans toutes les directions.

\begin{key}{Force de tension de surface}
Considérons un élément de surface $\dif S$ à l'interface entre un liquide et un gaz (par exemple, l'air). La tension de surface $\gamma$ agit tangentiellement à cette interface, et la force résultante par unité de longueur est $\gamma$. Pour un contour fermé de longueur $L$ délimitant une surface, la force totale due à la tension de surface est donc :
\begin{equation}
  \vec{F} = \gamma L \vec{t}
  \label{eq:tension_surface}
\end{equation}
où $\vec{t}$ est un vecteur unitaire tangent à la surface, perpendiculaire au contour $L$.

\end{key}

\section{Point de vue microscopique : accélération verticale des molécules}

Du point de vue microscopique, une molécule en surface est attirée par les molécules situées en dessous d’elle, mais elle n’a pas de voisine au-dessus pour la retenir. On en déduit qu’elle subit, en moyenne, une accélération dirigée vers le bas. Cependant, les chocs répulsifs avec les molécules voisines (analogues à des collisions élastiques) compensent cette attraction, de sorte que sa position moyenne reste stationnaire.

\section{Point de vue thermodynamique : tension de surface}

Une molécule dans la masse du fluide (le \emph{bulk}) est liée à ses voisines par un potentiel de Lenard-Jones, négatif. Une molécule en surface du fluide est liée à moins de molécules voisines. On en conclut donc que son \cindex{potentiel de liaison} est moins négatif qu'une molécule éloignée de la surface. Ainsi, pour une même masse de fluide, agrandir l'interface requiert de rompre des liaisons, ou encore d'augmenter l'énergie libre de la masse de fluide.

On note $\gamma$ le coefficient de \cindex{tension de surface}, défini comme le rapport entre la variation d’énergie libre $\dif H$ et la variation d’aire $\dif A$ de l’interface avec le vide :

\begin{displaymath}
  \gamma = \od{H}{A}
\end{displaymath}

Le travail nécessaire pour agrandir la surface est égal à la variation d’énergie libre $H$.

Pour une double interface vide -- liquide -- vide présentant un côté $L$, on a :
\begin{displaymath}
  \dif H = 2 \gamma L \dif x = F \dif x \Longrightarrow  F = 2 \gamma L.
\end{displaymath}

\section{Surface courbe}

\begin{figure}
  \begin{center}
    \input{Pdftex/07_tension_1.pdf_tex}
  \end{center}
  \caption{Bilan des forces exercées sur une parcelle de fluide en surface (générant une force nette à l'intérieur de la courbure).}\label{fig:07_tension_1}
\end{figure}

\begin{figure}
  \begin{center}
    \input{Pdftex/08_Young_Laplace.pdf_tex}
  \end{center}
  \caption{Illustration d'une interface liquide-gaz courbée : la différence de pression $\Delta p$ à travers une interface courbe est proportionnelle à la courbure locale et à la tension de surface $\gamma$.}\label{fig:08_Young_Laplace}
\end{figure}

Pour une interface courbe décrite par $z(x)$, la différence de pression $\Delta p$ entre les deux côtés de l’interface est donnée par la loi de Young-Laplace. En une dimension, cette relation s’écrit :

\begin{displaymath}
  \Delta p = \gamma \pd[2]{z(x)}{x}
\end{displaymath}

Cette équation montre que $\Delta p$ est proportionnelle à la courbure locale de l’interface (deuxième dérivée de $z$ par rapport à $x$) et à la tension de surface $\gamma$.


Si on considère une courbe un trois dimenions, c'est-à-dire le cas où $z$ devient une fonction $z(x,y)$, les tensions causées par la courbure selon les deux axes principaux s'additionnent. 
On montre alors que, si on définit l'axe $z$ comme étaint perpendiculaire à l'interface, et les axe $x$, $y$ selon les directions propres dela courbure, 

\begin{displaymath}
  \Delta p = \gamma (\pd[2]{z}{x} + \pd[2]{z}{y})
\end{displaymath}

C'est la loi de Young-Laplace. 

Elle est souvent présentée un peu autrement: 

\begin{studentinvestigation}{Calcul du rayon de courbure}
  Soit la courbe, dans le plan $x,z$, paramétrée de façon suivante: $x^2 + (z-R)^2  = R^2$.  C'est un cercle de rayon $R$, centrée en $(x,z)=(0,R)$. Ce cercle est tangent à la droite hozizontale $z=0$, au point $(0,0)$. Vérifiez que la partie inférieure du cercle répond à l'équation $z(x)=R-\sqrt{R-x^2}$. Ensuite calculez $\pd[2]{z}{x}$ au point de contact. Vous obtenez ..... $1/R$. 
\end{studentinvestigation}

Sachant en effet que la dérivée seconde est \cindex{rayon de courbure} du cercle tangeant à la courbe, on réécrit: 

\begin{key}{Loi de Young-Laplace}
  La différence de pression causée par la courbure d'une surface vaut
  \begin{equation}
  \Delta p = \gamma \left(\frac{1}{R_1} + \frac{1}{R_2}\right), 
    \label{eq:Young-Laplace}
  \end{equation}
où $R_1$ et $R_2$ sont les rayons de courbure principaux de la surface courbe. Le déficit de pression se trouve du côté convexe (la pression est plus élévée du côté concave, à l'intérieur de la surface). 
\end{key}

\begin{application}{Surpression à l'intérieur d'une goutte}
  Nous pouvouns nous convaincre de la cohérence de la théorie ci-dessous en calculant la surpression à l'intérieur d'une goutte de deux façon différence. La plus simple est d'utiliser la Loi de Young Laplace. Pour une goutte sphérique de rayon $R$, nous avons $\Delta p = \gamma \left(\frac{1}{R}+\frac{1}{R}\right) = \frac{2\gamma}{R}$. 

  Cependant nous pouvons également un bilan hydrostatique: la force exercée sur une surface plane coupant la goutte en deux ($\Delta p \pi R^2$) doit compenser très exactement la force, exercée dans l'autre sens, qui naît de la tension de surface, telle que définie en \eqref{eq:tension_surface}: $\gamma L = \gamma 2\pi R$. 
  \begin{center}
    \input{Pdftex/07_goutte.pdf_tex}
  \end{center}

\end{application}

Si on à affaire à une \cindex{bulle} plutôt qu'à une \cindex{goutte}, il faut tenir compte des tensions générées par les deux interfaces, air-eau, et eau-air par exemple. La surpression est donc doublée: à l'intérieur d'une goutte, la surpression est de $4\gamma/R$, où $\gamma$ est la tension de surface associée à l'intérface gaz-liquide. 

Nous venons d'éclaircir un mystère: sachant que la tension de surface de l'intérface air-eau est d'environ $72$\ mN/m, et que la tension de surface à l'intérface air-eau savonneuse est plutôt de l'ordre de $25$\ nM/m, nous voyons pourquoi les bulles d'eau savonneuse sont beaucoup plus stables que celle d'eau. 

\section{Interfaces triples: Young - Laplace}



\section{Loi de Jurin}


\todo[inline]{A faire}
